\subsection{Пространственная и временная дискретизация}
\label{subsec:coverage_discretization}

Для численной оценки область $\mathcal{R}$ заменяется сеткой H3 заданного
разрешения. Обозначим через
\begin{equation}
    \mathcal{Q}_{\rho}=\{Q_1,\ldots,Q_{N_{\rho}}\}
    \label{eq:h3_point_set}
\end{equation}
множество центров H3-сот разрешения $\rho$, выбранных для рассматриваемой
области. Тогда численная оценка показателя покрытия
из~\eqref{eq:coverage_metric} имеет вид
\begin{equation}
    \widehat{\mu}_{\rho}(t,\mathbf{x})=
    \frac{\left|\mathcal{Q}_{\rho}\cap
        \mathcal{C}(t,\mathbf{x})\right|}{N_{\rho}}.
    \label{eq:discrete_coverage_metric}
\end{equation}
Значение $\widehat{\mu}_{\rho}(t,\mathbf{x})=1$ означает, что в момент $t$
покрыты центры всех выбранных H3-сот. Результат такой оценки зависит от
разрешения $\rho$: его увеличение повышает пространственную детальность
проверки, но требует больших вычислительных затрат.

Для проверки покрытия во времени непрерывный интервал $\mathcal{I}$
заменяется конечной расчётной сеткой
\begin{equation}
    \mathcal{T}=\{t_0,t_1,\ldots,t_M\}\subset\mathcal{I}.
    \label{eq:time_grid}
\end{equation}
Начальная эпоха, границы интервала, модель движения и плотность сетки
$\mathcal{T}$ задаются до начала численного решения. Временной шаг выбирается
с учётом характерных угловых масштабов построения, включая расстояние между
соседними КА и ширину полосы покрытия. В результате условие покрытия из
задачи~\eqref{eq:optimization_problem} проверяется в дискретном виде:
\begin{equation}
    \min_{t\in\mathcal{T}}
        \widehat{\mu}_{\rho}(t,\mathbf{x})=1.
    \label{eq:discrete_time_coverage}
\end{equation}
